% !TEX root = ../main.tex
\chapter{Idriftsættelse og test}
\label{ch:idrift}

Dette kapitel dokumenterer den software- og styringsmæssige idriftsættelse
af armen efter den mekaniske og elektriske integration. Kapitlet beskriver
i hvilken konfiguration robotten faktisk køres, hvilke værktøjer der er
udviklet for at kunne måle armens opførsel, hvad disse værktøjer faktisk
måler -- og lige så vigtigt: hvad de \emph{ikke} måler -- samt en række
konkrete problemer, der er stødt på undervejs, og hvordan de er løst.

\section{Software-opsætning}

Idriftsættelsen tager udgangspunkt i Energinets eksisterende ROS\,2-workspace
(\enquote{energirobotter-ros-workspace}), der ligger på Jetson Orin Nano under
ROS\,2 Humble. For at kunne arbejde isoleret med én arm uden at indlæse
konfigurationer for hænder, hoved og højre arm, indlæses kun den relevante
\enquote{servo\_arm\_left\_params.json}-fil ved opstart af
\enquote{wattson\_servo\_manager\_node}. Selve opstarten sker via to launch-filer:
\enquote{servos.launch.py} starter manager-noden, og
\enquote{slider\_control.launch.py} starter en GUI med en skyder pr.~led, der
gør det muligt at bevæge hvert led uafhængigt.

\section{PID-konfiguration og kontrolarkitektur}

Hver ST3215-servo har en intern positionsregulator. PID-værdierne
pr. servo står i JSON-filen og er anvendt uændret fra Energinets
konfiguration. Manager-noden videresender målpositioner til hver
servo, og selve reguleringssløjfen lukkes inde i den enkelte ST3215.

ROS-laget regulerer ikke selv og modtager ingen tilbagemelding fra
servoernes faktiske position. Som test blev \enquote{feedback\_enabled}
sat til \enquote{true} i både JSON-konfig og driverkode for at
undersøge, om feedback-vejen kunne aktiveres uden yderligere arbejde.
Resultatet var, at alle arm-led begyndte at oscillere, også uden
eksterne kommandoer. Open-loop-konfigurationen blev derfor
reetableret, og feedback-vejen blev i stedet brugt diagnostisk
(afsnit~\ref{sec:hvad-maaler-vi}).

\section{Kalibrering af led}

Inden armen kan styres meningsfuldt, skal hvert led kalibreres. For hver
servo indeholder JSON-konfigurationen tre værdier, der definerer leddets
arbejdsrum:

\begin{itemize}
    \item \enquote{angle\_software\_min} og \enquote{angle\_software\_max} --
          de software-grænser, manager-noden klipper kommandoer til, så
          hverken servo eller mekanik ødelægges af en ude-af-grænse-kommando.
    \item \enquote{default\_position} -- den vinkel, leddet sættes til ved opstart.
    \item \enquote{gear\_ratio} -- omsætningsforholdet mellem servo-aksel
          og fysisk led, så ROS-niveauet kan tænke i led-vinkler
          frem for servo-aksel-vinkler.
\end{itemize}

Kalibreringen er sket praktisk: armen bevæges manuelt med
slider-GUI'en, og grænserne sættes til de positioner, hvor leddet
fysisk når sit endestop uden at trække i mekanikken. Nul-positionen er
sat til en geometrisk neutral pose, hvor armen hænger lodret langs
torsoens side. Værdierne for venstre arm er dokumenteret i bilaget
(se \emph{Bilag, Tabel~1}).

\section{Tilføjelser: commissioning-værktøjer}
\label{sec:commissioning-vaerktoejer}

Det oprindelige system har ingen indbyggede værktøjer til at måle eller
dokumentere armens dynamiske opførsel. Til brug for idriftsættelsen er
der derfor udviklet en ny ROS-pakke, \enquote{arm\_commissioning}, der
indeholder følgende noder:

\begin{description}
    \item[\enquote{step\_response\_node}] Kommanderer en step-input på en
        valgt led, logger den feedback-vinkel, manager-noden publicerer,
        og beregner klassiske step-respons-mål: rise time (10\,\%--90\,\%),
        overshoot, settling time (med konfigurerbar tolerance) og
        steady-state-fejl. Output: CSV-rådata og en PNG-graf pr.~kørsel.
    \item[\enquote{repeatability\_node}] Sender armen frem og tilbage
        mellem to konfigurerede positioner et antal cyklusser og
        registrerer slut-positionen i hvert hold-vindue. Output: histogram
        for de to målte pose-populationer plus mean, std og max-afvigelse.
    \item[\enquote{launcher\_gui\_node}] Et Tk-baseret betjeningspanel, der
        samler de typiske demo-services (DHCP, kameraets ROS-launch,
        VR-broer, servo-manager, animationsafspilning) bag start/stop-knapper
        med statusindikatorer og samlede logfaner. Begrundelse: en
        eksamensdemo, hvor 5--6 terminalvinduer skal åbnes i en præcis
        rækkefølge, er fragil; et samlet panel reducerer antallet af
        fejlmuligheder.
\end{description}

For at \enquote{step\_response\_node} og \enquote{repeatability\_node} kunne
hente positionsdata, blev der desuden tilføjet en feedback-publisher til
\enquote{wattson\_servo\_manager\_node}: et nyt topic
\enquote{/joint\_states\_feedback}, der publicerer den positionsværdi,
manager-noden kender for hver servo. Dette topic er den eneste nye
afhængighed, commissioning-værktøjerne tilfører til det eksisterende
system.

\section{Funktionstest: step-response pr.~led}

Step-response køres på hvert af de syv venstre arm-led med samme
parametre: 10\textdegree{} step fra nul, 12\,s logging, 50\,Hz
samplerate, settling-tolerance 10\,\% af step-amplituden.

% TODO ALEX (efter lab): indføj de målte værdier i tabellen
% (Bilag, Tabel~2) og det repræsentative plot (Bilag, Figur~1).
% Hvis tallene afslører større afvigelser mellem leddene, så
% diskutér dem i kapitel~\ref{ch:diskussion}.
Resultatet pr.~led sammenfattes i tabellen i bilaget
(se \emph{Bilag, Tabel~2}). Et repræsentativt forløb for
\enquote{joint\_left\_shoulder\_pitch} vises i \emph{Bilag, Figur~1}:
den kommanderede vinkel som step-funktion, den loggede
feedback-vinkel og de beregnede mål annoteret på plottet.

% TODO efter lab lørdag: udfyld med de faktiske målte værdier.
Afvigelser mellem leddene diskuteres i kapitel~\ref{ch:diskussion}.

\section{Hvad værktøjerne måler -- og ikke måler}
\label{sec:hvad-maaler-vi}

Den feedback-vinkel, \enquote{step\_response\_node} logger, kommer fra
ROS-topic'et \enquote{/joint\_states\_feedback}. Det topic publiceres
af manager-noden og indeholder i den nuværende konfiguration den
\emph{senest kommanderede} vinkel for hvert led, ikke en målt
position fra servoens encoder. Det skyldes en tre-niveau gating på
flaget \enquote{feedback\_enabled}, der i den oprindelige konfiguration
står til \enquote{false}:

\begin{enumerate}
    \item I JSON-konfigurationen for venstre arm er
          \enquote{feedback\_enabled} sat til \enquote{false} for alle
          syv servoer.
    \item Driveren \enquote{DriverWaveshare} har samme værdi som
          instans-default, og manager-noden overskriver den ikke.
    \item I \enquote{ServoControl.compute\_command} medfører
          \enquote{feedback\_enabled = false}, at den cachede position
          \enquote{self.angle} sættes lig den senest kommanderede vinkel
          frem for en målt værdi.
\end{enumerate}

Konsekvensen for commissioning-værktøjerne er, at de tal, de
udregner, ikke karakteriserer ST3215-servoens interne regulator. Det
de måler, er kommandoernes vej gennem systemet -- fra en ny target
publiceres på \enquote{/joint\_states}, til den cachede position
afspejler den nye target. Det er en nyttig størrelse, fordi den
dækker latency i ROS-laget, men det er \emph{ikke} en
servo-karakteristik.

Som beskrevet under PID-konfigurationen blev det afprøvet at slå
\enquote{feedback\_enabled} til; armen begyndte at oscillere, og
open-loop-konfigurationen blev derfor bibeholdt. Step-respons-tabellen
i bilaget skal læses i lyset af dette: tallene er gyldige som
sammenligningsgrundlag mellem leddene og som karakteristik af
command-pipelinen, men de er ikke et udtryk for ST3215-regulatorens
dynamik.

\section{Andre integrationsrettelser}

Ud over de fund, der direkte vedrører funktionstesten, blev der under
idriftsættelsen rettet en række mindre integrationsfejl. De vigtigste
er listet nedenfor; detaljer ligger i kildekodens commit-historik og
bilag.

\begin{description}
    \item[Stabil portbinding for ESP32-bokse.] Tidligere blev de tre
        USB-tilsluttede mikrocontrollere adresseret som
        \enquote{/dev/ttyUSB0--2}, hvilket gjorde tildelingen afhængig af
        tilslutningsrækkefølge. Manager-noden adresserer dem nu via
        \enquote{/dev/serial/by-path/}, så hver fysisk port på
        Jetson-USB-hubben permanent svarer til den samme logiske funktion.
        Detaljer i kapitel~\ref{ch:el}.
    \item[Race-tilstand i kinematik-nodens opstart.]
        \enquote{elrik\_kdl\_kinematics\_node} crashede sporadisk under
        opstart med en \enquote{AttributeError}, fordi en periodisk timer
        blev oprettet, før node-instansens egne felter var initialiseret.
        Fixet ved at flytte \enquote{create\_timer}-kaldet til slutningen
        af \enquote{\_\_init\_\_}, så timeren først aktiveres, når
        nodens tilstand er klar.
    \item[Launcher-GUI bugs under hardware-rehearsal.]
        Det samlede betjeningspanel blev iterativt rettet for en række
        problemer, der først kom frem under faktisk demo-brug. De fire
        vigtigste rettelser beskrives nedenfor.

        \paragraph{Adgangskode-prompts ved fjernlogin.}
        Launches af ROS-noder på Jetson sker fra GUI'en via en
        krypteret fjern\-forbindelse. Oprindeligt skulle operatøren
        indtaste password ved hvert kald, hvilket gjorde et
        sammenhængende demo-flow uigennemførligt. Problemet blev løst
        med to mekanismer. Først blev der lavet en lille hjælpe-proces,
        som leverer adgangskoden gennem et grafisk dialog\-vindue, så
        fjernforbindelsen ikke længere kræver, at GUI'en har en
        terminal at spørge i. Dernæst blev der konfigureret deling af
        forbindelsen, så det første login etablerer en
        baggrunds-forbindelse, som efterfølgende kald genbruger.
        Resultatet er, at password kun spørges én gang pr.~service og
        ikke én gang pr.~ROS-launch. En mere permanent løsning ville
        være login med kryptografisk nøgle frem for adgangskode, men
        det er fravalgt, da det ville kræve ændringer i Jetsons
        konfiguration uden for projektets afgrænsning.

        \paragraph{Langsom respons på stop-knappen.}
        Den oprindelige stop-knap i GUI'en reagerede med 30--60
        sekunders forsinkelse, hvilket gjorde nødstop praktisk
        ubrugeligt. Årsagen var, at stop-signalet ikke nåede helt ned
        til de ROS-noder, der faktisk kørte: signalet blev fanget af
        en mellemliggende skal-proces, og selve noden levede videre
        som en frikoblet baggrunds-proces. Rettelsen bestod i tre
        dele. For det første blev kommandoen om\-skrevet, så
        ROS-noden overtager skal-processen direkte i stedet for at
        køre som dens barn. For det andet startes hver service nu i
        sin egen proces-gruppe, så stop-knappen kan signalere hele
        gruppen på én gang og dermed også fange middleware- og
        controller-processer, der eventuelt er startet sammen med
        noden. For det tredje sendes først et blødt afbrydelses-signal,
        som lader noden rydde op, og hvis processen mod forventning
        stadig kører efter en kort frist, sendes et hårdere
        terminerings-signal som fallback. Tilsammen bringer
        ændringerne respons\-tiden under et sekund.

        \paragraph{Manglende synkronisering mellem laptop og Jetson.}
        Under en refaktorering blev kildekoden til en central node
        på laptop-siden uforvarende beskadiget, men på Jetson kørte
        systemet tilsyneladende videre uden fejl. Forklaringen var,
        at Python på Jetson havde en cached udgave af samme modul
        liggende fra en tidligere build og foretrak den frem for at
        fejle på den nye kilde. Diskrepansen blev først opdaget, da
        Jetson blev rebuildt fra grunden og pludselig fejlede ved
        opstart. Den varige rettelse blev derfor proces-mæssig og
        ikke kode-mæssig: før en kritisk demo skal Jetson altid
        bygges fra rene mapper, så det testede svarer til det
        kørende.

        \paragraph{Øvrige mindre rettelser.}
        Ved luk af GUI'en kunne et stop-kald ramme en system-tjeneste
        ejet af root og resultere i en rettighedsfejl, så vinduet
        hang. Det blev rettet ved at omslutte alle service-stop i
        fejl\-håndtering og foretage netop dette stop med forhøjede
        rettigheder. Dernæst blev et eksperiment med at hæve
        kontrol-frekvensen fra 10 til 50~Hz for jævnere bevægelse
        rullet tilbage igen. Robottens 14 servoer deler en fælles
        seriebus, og målingerne pr.~servo er begrænset af bussens
        gennemløb, ikke af CPU'en. En cyklustid på 20~ms svarende
        til 50~Hz er derfor fysisk umulig på den eksisterende bus,
        og forsøget gav i praksis ustabile bevægelser uden gevinst.
    \end{description}

\section{Resultat}

Ved afleveringstidspunktet kan armen styres positionsmæssigt led for led
med slider-GUI'en og kan afspille præindspillede animations-sekvenser
fra commissioning-pakkens animations-fane. Pålideligheden af de
underliggende ROS-topics er verificeret via \enquote{ros2 topic hz}.

% TODO efter lab lørdag: bekræft at step-respons er kørt for alle
% syv led, eller juster sætningen herunder hvis kun en del er kørt.
Step-respons køres på alle syv venstre arm-led og dokumenteres i bilaget.

De væsentligste ikke-opnåede mål er en kalibreret closed-loop
positionsregulering på ROS-niveau (fravalgt, jf.~ovenfor) og en
funktionsdygtig per-led-effektmåling; sidstnævnte er påbegyndt, men
ikke leveret, og er diskuteret i kapitel~\ref{ch:diskussion}.
